\documentclass[11pt]{article}
\usepackage{ctex}
\usepackage[a4paper, margin=2.5cm]{geometry}
\usepackage{booktabs}
\usepackage{array}
\usepackage{enumitem}
\usepackage{xcolor}
\usepackage{hyperref}
\usepackage{fancyhdr}
\usepackage{titlesec}
\usepackage{verbatim}

\setCJKmainfont{FandolSong}
\setCJKsansfont{FandolHei}

\hypersetup{
    colorlinks=true,
    linkcolor=blue,
    urlcolor=blue,
    pdftitle={IAI2O 备赛培训 - 开营第一课 项目说明},
    pdfauthor={张浩然}
}

\titleformat{\section}{\Large\bfseries}{\thesection}{1em}{}
\titleformat{\subsection}{\large\bfseries}{\thesubsection}{1em}{}
\titleformat{\subsubsection}{\normalsize\bfseries}{\thesubsubsection}{1em}{}

\setlength{\headheight}{14pt}
\pagestyle{fancy}
\fancyhf{}
\fancyhead[L]{IAI\textsuperscript{2}O 备赛培训}
\fancyhead[R]{开营第一课}
\fancyfoot[C]{\thepage}

\title{\textbf{IAI\textsuperscript{2}O 备赛培训 · 开营第一课 项目说明}}
\author{主讲：张浩然}
\date{2026 年 06 月 12 日}

\begin{document}

\maketitle

\begin{quote}
课程：IAI\textsuperscript{2}O 备赛培训 · 开营第一课《认识 / 看清 / 出发》
\end{quote}

\section{核心理念}

AI 已将知识与技术的门槛几乎推平，但"人人都能用"不等于"人人用得一样好"。

会用 AI 的人能将效率提升 10--20 倍，而不会用的人只能获得 2--3 倍的提升。
\textbf{差距不在有没有 AI，而在会不会问、会不会判断、能不能把它串成一件事。}

\subsection*{AI 不再是瓶颈——你才是}

决定 AI 使用上限的四件事（与会不会写代码无关）：

\begin{enumerate}
    \item \textbf{问对问题} — 问题提得准，AI 才能给出有价值的回答
    \item \textbf{分得出高下} — 在多个方案中判断哪个更优，这是品味
    \item \textbf{定方向、立标准} — 知道目标在哪、什么叫"做对了"
    \item \textbf{拆解与编排} — 把大目标拆成 AI 能执行的步骤并串联起来
\end{enumerate}

\section{工具库}

\begin{table}[h]
\centering
\begin{tabular}{>{\bfseries}p{3.2cm} p{2.8cm} p{8cm}}
\toprule
工具 & 定位 & 核心用途 \\
\midrule
Claude Code & 主力 AI Agent & 自动读文件、写代码、跑程序、查报错，兼具导师和讨论伙伴角色 \\
GitHub & 项目中枢 & 版本管理、全程留痕、多人协作、任务看板 \\
LaTeX / Overleaf & 论文排版 & 在线编写并编译出专业格式 PDF，满足比赛交付要求 \\
AutoDL & 云端算力 & 按小时租用 GPU，训练模型无需本地显卡 \\
\bottomrule
\end{tabular}
\end{table}

\section{比赛目标}

\textbf{IAI\textsuperscript{2}O（国际人工智能创新奥林匹克）}

\begin{itemize}
    \item 面向全球 13--18 岁中学生
    \item 2026 年 9 月在 MIT 线下总决选
    \item 评委来自哈佛 / MIT / 伯克利，院士级科学家委员会
    \item 我们参赛方向：\textbf{AI4Sci --- 用 AI 做考古遗址发现与遗产保护}
\end{itemize}

\subsection*{比赛交付物}

\begin{enumerate}
    \item 英文研究论文（LaTeX 排版）
    \item 研究日志 / 可复现记录（GitHub 全程留痕）
    \item 四象限图（形态置信度 $\times$ 受威胁程度）
    \item 展板 + 英文答辩
\end{enumerate}

\section{项目方案：AI 辅助考古遗址发现与保护}

\subsection*{问题背景}

全球遥感影像每天持续增长，但能判读考古遗迹的专家严重不足：

\begin{itemize}
    \item \textbf{发现端}：该被发现的遗址无人处理
    \item \textbf{守护端}：盗掘破坏痕迹无人持续监测
\end{itemize}

\textbf{瓶颈}：专家注意力的稀缺。

\subsection*{方法论：望远镜逻辑}

\begin{table}[h]
\centering
\begin{tabular}{>{\bfseries}p{3.5cm} p{10cm}}
\toprule
步骤 & 内容 \\
\midrule
Step 1 · 造仪器 & 搭建从影像自动识别疑似遗迹的 AI 工具链 \\
Step 2 · 标定仪器 & 在有专家标注的 Arran 区域（700+ 条记录）验证准确性 \\
Step 3 · 用仪器观测 & 对准真实目标区域，输出可信结果 \\
\bottomrule
\end{tabular}
\end{table}

\subsection*{AI 工具链流程}

\begin{verbatim}
影像（LiDAR / 卫星影像）
    |
    v
YOLO 检测 -- 框出疑似遗迹
    |
    v
LLM 裁决 -- 按形态学判据逐个判断"是 / 不是"
    |
    v
去重 -- 比对国家遗址库（Canmore 31 万条）
    |
    v
证据档案 -- 自动生成候选清单
\end{verbatim}

\subsection*{两条观测线}

\begin{itemize}
    \item \textbf{发现线 · 苏格兰 Kintyre}：对从未系统排查过的半岛做普查，汇出新候选清单提交官方
    \item \textbf{守护线 · 阿富汗}：对已知遗址的多年卫星影像逐月比对，生成盗掘监测报告
\end{itemize}

\section{本次作业目标}

\begin{quote}
把今天讲的工具自己串起来跑通一遍，体验"会指挥 AI"的完整流程。
\end{quote}

\begin{enumerate}
    \item 配置 Claude Code 与 GitHub 权限
    \item 用 Claude 将课件总结为项目说明（即本文件）
    \item 用 Claude Code 新建 GitHub 仓库并上传
    \item 用 Claude Code 将总结用 LaTeX 排版、编译成 PDF，一并上传
\end{enumerate}

\begin{center}
\textbf{你只负责下指令、做判断，其余交给 Claude Code。}
\end{center}

\end{document}
